% Owner: Fatima Malik / Member B.
% Local font encoding restores the template's Times regular/bold faces under XeTeX.
\begingroup
\renewcommand{\encodingdefault}{T1}
\fontencoding{T1}\selectfont
\chapter{Literature Review / Related Work}
Clinical decision support, medical language models, imaging analysis and retrieval systems each address a part of the problem that Zeest brings together. Existing clinical tools can provide references, differential diagnosis suggestions or patient record support, but they do not by themselves establish a longitudinal, physician-controlled, multimodal reasoning workflow. Recent language models show strong performance on medical question answering, while biomedical language models improve representation of clinical and scientific text. Imaging models and chest radiograph datasets provide validated groundwork for supported diagnostic inputs. Retrieval-augmented generation and medical retrieval methods show how external evidence can be supplied to a generator, and agentic reasoning work motivates stepwise decomposition. The related work is therefore reviewed by capability area rather than as a single comparison list. This structure makes the gap clear: Zeest is not proposed as a new foundation model or a new imaging classifier, but as an integrated clinical workflow that connects memory, multimodal findings, cited evidence, critic review and explicit physician approval.

\section{Definitions, Acronyms, and Abbreviations}
\textbf{CDSS:} Clinical Decision Support System.

\textbf{EMR:} Electronic Medical Record.

\textbf{EHR:} Electronic Health Record.

\textbf{LLM:} Large Language Model.

\textbf{NLP:} Natural Language Processing.

\textbf{RAG:} Retrieval-Augmented Generation.

\textbf{AI:} Artificial Intelligence.

\textbf{CXR:} Chest X-ray.

\textbf{MIMIC-IV:} Medical Information Mart for Intensive Care IV.

\textbf{MIMIC-CXR:} Medical Information Mart for Intensive Care Chest X-Ray.

\textbf{BioGPT:} Biomedical Generative Pre-trained Transformer.

\textbf{ClinicalBERT:} Clinical Bidirectional Encoder Representations from Transformers.

\textbf{MedCPT:} Medical Contrastive Pre-trained Transformer.

\textbf{ADK:} Agent Development Kit.

\textbf{Ablation Study:} A planned comparison in which selected components are removed or added to examine their contribution.

\textbf{Longitudinal Memory:} Patient-specific clinical context retained and retrieved across encounters.

\textbf{Multimodal Input:} Clinical information represented through more than one type of data, such as notes, imaging, laboratory information and medication context.

\section{Detailed Literature Review} 
This section reviews systems, models, datasets and methods related to Zeest. The topic is divided into clinical decision support, medical language models, imaging and multimodal learning, retrieval-based evidence support, agentic orchestration and clinical data resources.

The purpose of this review is to identify what existing work already supports and where Zeest remains different. Thus, the literature is discussed according to the capability it contributes to the proposed physician-controlled, longitudinal and multimodal clinical workflow.

\subsection{Clinical Decision Support and Electronic Records}
\subsubsection{Summary of the research item}
Isabel DDx represents a differential diagnosis support application that helps clinicians consider possible diagnoses from entered clinical features.\cite{zeestB:isabel} Its strength is that it addresses a real clinical reasoning task and can broaden the physician's diagnostic consideration. Its limitation for Zeest is that it is not described as a longitudinal patient-memory system with multimodal agents, audit-controlled approval and an ablation framework. UpToDate provides evidence-based clinical reference material for clinicians and is useful as a point of care knowledge resource.\cite{zeestB:uptodate} It supports evidence access, but it does not organise patient-specific memory or produce a physician-reviewed draft through specialist agents. OpenMRS demonstrates how an open source EMR can support patient record management, especially in resource-constrained settings.\cite{zeestB:openmrs} Its relevance is strongest for patient and encounter organisation, while its limitation is that it is a record platform rather than an agentic reasoning and evidence-review workflow.

\subsubsection{Critical analysis of the research item}
The reviewed systems are strong in their individual areas. Isabel DDx supports differential diagnosis exploration, UpToDate gives clinicians access to curated medical reference content and OpenMRS shows how patient records can be organised in an open source EMR. Their common weakness is that each addresses only part of Zeest's target workflow. Differential diagnosis support does not automatically preserve encounter-specific memory, a reference system does not manage patient context, and an EMR does not by itself generate evidence-grounded drafts through specialist agents. These systems also do not define revision-specific approval, modification, rejection and audit handling. Their limitations show why Zeest needs a combined method rather than one existing clinical tool.

\subsubsection{Relationship to the proposed research work}
These systems show that Zeest should not treat clinical decision support as only a text-generation task. Differential diagnosis applications, medical references and EMR platforms each supply useful ideas for the project: possible diagnosis review, cited clinical knowledge and structured patient records. However, the proposal problem concerns the connection between these ideas. A physician needs to view the current consultation, retrieve relevant longitudinal history, inspect specialist findings and approve or reject the resulting draft. Existing reference or EMR systems do not automatically answer how a generated recommendation should be linked to a specific encounter, supporting evidence and audit action. Chapter~4 therefore specifies patient selection, encounter history, consultation note correction, evidence inspection and revision-specific approval as separate behaviours. The literature supports those boundaries by showing that information access and decision support are different responsibilities. Zeest's contribution is the proposed integration of those responsibilities under physician control, not replacement of existing clinical judgement.

\subsection{Medical and Biomedical Language Models}
\subsubsection{Summary of the research item}
Med-PaLM 2 demonstrates that large language models can achieve strong performance on medical question answering benchmarks after medical alignment and evaluation.\cite{zeestB:medpalm2} GPT-4 medical challenge work also shows broad medical reasoning capability across examination-style tasks, while documenting safety and reliability concerns that matter in clinical use.\cite{zeestA:noriMedical} These studies motivate the use of a powerful language model within Zeest, but their task setting is different from a longitudinal clinical workflow. Benchmark questions do not require storing a patient's encounter history, associating outputs with a physician action or preserving the difference between a draft and an approved clinical decision. The literature therefore supports cautious use of language models as reasoning components. It does not support autonomous diagnosis, unreviewed treatment or the assumption that benchmark performance proves patient-level reliability. Zeest addresses this limitation by requiring retrieved evidence, critic review and explicit physician approval before any generated content becomes a recorded clinical decision.

\subsubsection{Critical analysis of the research item}
ClinicalBERT and BioGPT show another important direction: adapting language models to biomedical and clinical text. ClinicalBERT provides publicly available contextual embeddings trained on clinical notes, making it relevant to patient-record representations and clinical NLP tasks.\cite{zeestB:clinicalbert} BioGPT applies generative pre-training to biomedical text generation and mining, showing that domain-specific corpora can improve biomedical language processing.\cite{zeestB:biogpt} Their strength is representation of specialised language, terminology and document patterns. Their limitation for Zeest is that they do not by themselves provide a complete physician workflow, multimodal processing path or approval-controlled memory update. They also do not solve the problem of grounding a draft in patient-specific evidence. For Zeest, these works support the general choice to use medical-language capability and semantic representations, but the actual embedding model and LLM remain unresolved. The methodology must therefore treat model selection as a dependency and avoid claiming that any specific biomedical model has already been integrated.

\subsubsection{Relationship to the proposed research work}
These language-model studies are related to Zeest because the proposed system requires clinical text understanding, structured note use, reasoning over case context and generation of a reviewable draft. Their results support the decision to include language-model capability, but they also explain why Zeest requires additional controls. The proposed work adds longitudinal patient memory, cited evidence retrieval, critic review and physician approval around the generative component. In this way, Zeest uses medical language modelling as one part of a controlled clinical workflow rather than treating a model answer as an accepted clinical decision.

\subsection{Medical Imaging and Multimodal Learning}
\subsubsection{Summary of the research item}
CheXNet used deep learning for pneumonia detection on chest radiographs and is frequently cited as evidence that imaging models can reach strong performance on a narrow radiology task.\cite{zeestB:chexnet} CheXpert then provided a large chest radiograph dataset with uncertainty labels and an expert comparison benchmark across fourteen observations.\cite{zeestB:chexpert} These works are relevant because the Zeest proposal intends to integrate an existing validated imaging classification model rather than train a new diagnostic model from scratch. Their limitation is equally important: a chest X-ray classifier produces findings for a supported imaging task, not a complete clinical decision. It does not know the patient's longitudinal history, current consultation, medication context or physician approval state. Therefore, Chapter~5 must place imaging analysis as one input to the clinical case rather than the central method. The supported clinical domain depends on a dataset and validated model that can be used under the project constraints.

\subsubsection{Critical analysis of the research item}
BioViL-T extends the imaging discussion by learning from temporal structure in biomedical vision-language processing.\cite{zeestB:biovilt} It is relevant to Zeest because the proposal is concerned with more than a single isolated input; the system must reason over patient information across time and modalities. However, a vision-language model still differs from the proposed Zeest workflow. Zeest combines consultation notes, history retrieval, imaging outputs, laboratory information, medication context, evidence retrieval and critic review. A multimodal representation alone does not establish the audit trail or decision boundary required by Chapter~4. The review of imaging and vision-language research therefore supports the need for careful integration. Imaging findings should retain their source and be displayed as specialist findings. They should not be silently merged into a generated conclusion without indicating how they relate to the current encounter. This distinction also motivates the planned text-only versus multimodal comparison in Chapter~5.

\subsubsection{Relationship to the proposed research work}
The imaging literature directly supports Zeest's plan to integrate an existing validated model instead of training a diagnostic model from scratch. It also supports the need for a multimodal methodology in which imaging findings are considered alongside notes, history, laboratory information and medication context. Zeest differs from isolated imaging systems by treating image analysis as a specialist contribution to a physician-reviewed case. The final draft remains subject to evidence retrieval, critic review and explicit physician approval, so imaging output is not treated as an independent clinical decision.

\subsection{Retrieval, Evidence Grounding and Medical Search}
\subsubsection{Summary of the research item}
Retrieval-augmented generation connects a generator with retrieved documents so that generated answers can use external knowledge.\cite{zeestA:rag} This is central to Zeest because a clinical draft should be accompanied by evidence that the physician can inspect. General RAG, however, does not guarantee that retrieved passages are clinically appropriate, patient-specific or sufficient for a diagnosis. MedRAG examines retrieval-augmented generation in medical question answering and provides a closer medical context for evaluating retrieval and generation together.\cite{zeestB:medrag} Its relevance lies in showing that retrieval choices can affect medical answer quality. Its limitation is that medical question answering still differs from a physician-facing encounter workflow with patient memory and approval actions. MedCPT contributes contrastive pre-training for biomedical retrieval using PubMed search logs, supporting the idea that medical retrieval benefits from domain-specific search representations.\cite{zeestB:medcpt} Zeest can use these ideas only after selecting appropriate evidence sources and retrieval components.

\subsubsection{Critical analysis of the research item}
The strength of RAG and medical retrieval work is that it gives a practical way to connect generated text with external sources. MedRAG and MedCPT are especially relevant because they move retrieval closer to biomedical and medical question answering. The limitation is that retrieval does not automatically make an answer correct. Retrieved passages may be incomplete, irrelevant or poorly matched to the patient case. A generator may also overstate what the retrieved evidence supports. These weaknesses are important for Zeest because a clinical draft must be reviewable by a physician. The system must therefore keep evidence visible, cite sources and allow critic review instead of hiding retrieval inside the generated answer.

\subsubsection{Relationship to the proposed research work}
The evidence literature supports two design principles for the proposed system. First, evidence must remain visible and citable rather than disappearing into a generated answer. Second, retrieval quality is a variable to be evaluated, not an assumed property of the architecture. Chapter~4 reflects these principles through requirements for guideline retrieval, citation display and supporting passages. Chapter~5 extends them into a proposed evidence-retrieval procedure in which patient context and external guideline search are separated before the clinical draft is produced. The physician should be able to inspect whether a cited passage actually supports the draft. The Safety and Critic Agent should also review unsupported or contradictory content before presentation. This approach is stricter than ordinary RAG demonstrations because the final workflow concerns a clinical draft and an audit record. The literature therefore motivates RAG use, while the project constraints define the additional safeguards needed for Zeest.

\subsection{Agentic Reasoning and Orchestration}
\subsubsection{Summary of the research item}
ReAct combines reasoning traces with actions, showing how a language model can alternate between thinking and interacting with external tools.\cite{zeestB:react} Its relevance to Zeest is conceptual: clinical support can benefit from separating reasoning steps and tool use instead of relying on a single monolithic generation. LangGraph and Google ADK are named in the proposal as candidate orchestration technologies for agentic workflows.\cite{zeestB:langgraph,zeestB:googleadk} They are relevant because Zeest includes specialised Patient Memory, Clinical Case, Imaging, Laboratory, Medication Safety, Evidence and RAG, Clinical Reasoning, and Safety and Critic agents. However, the exact division between LangGraph and ADK remains unresolved in the repository rules and Kashaf's handoff. The literature and tool documentation therefore support the possibility of orchestration but do not settle Zeest's final implementation structure. Chapter~5 should describe the processing responsibilities and information flow without inventing a completed orchestration design.

\subsubsection{Critical analysis of the research item}
Agentic systems also introduce risks that are especially important in a clinical report. More components can improve separation of responsibility, but they can also create ambiguous handoffs, repeated context, failed tool calls or unsupported conclusions. Zeest's method must therefore define what each stage receives and returns. The Patient Memory Agent retrieves relevant history, the Clinical Case Agent organises the case, specialist agents produce typed findings, the Evidence and RAG Agent retrieves sources, the Clinical Reasoning Agent drafts the recommendation and the Safety and Critic Agent reviews it. The physician then approves, modifies or rejects the draft. This sequence is proposed, not executed. Its value will be examined through the planned ablation study rather than assumed in advance. ReAct and orchestration frameworks motivate the decomposition, while Zeest adds the patient-memory, evidence-inspection and physician-approval controls required for the medical setting.

\subsubsection{Relationship to the proposed research work}
This work supports Zeest's decision to divide responsibilities across specialised agents rather than using one prompt for the whole clinical task. The relationship is strongest in the idea that reasoning can be coordinated with tool use, retrieval and intermediate checks. Zeest extends that idea into a medical copilot workflow by adding patient memory, specialist clinical inputs, critic review and physician-controlled decision recording. The proposed research will compare this full agentic workflow against simpler variants so that the value of the decomposition can be assessed instead of assumed.

\subsection{Clinical Datasets and Synthetic Cases}
\subsubsection{Summary of the research item}
MIMIC-IV provides de-identified electronic health record data and is relevant to Zeest's longitudinal-memory objective because it contains patient and encounter-level clinical information.\cite{zeestA:mimicIV} MIMIC-CXR provides a de-identified chest radiograph dataset with associated reports, making it relevant to imaging and multimodal case preparation.\cite{zeestB:mimiccxr} These datasets support the idea that de-identified clinical data can be used for research, but their existence does not by itself grant access, determine the supported disease domain or authorize every planned use. Synthea offers a method and software mechanism for generating synthetic patients.\cite{zeestB:synthea} Its relevance is that synthetic cases can help develop workflows when real clinical data access is constrained. The limitation is that synthetic cases may not represent the full clinical complexity of real patients. Zeest's methodology therefore distinguishes selected candidates, access status, preprocessing and final evaluation resources instead of treating all proposal-listed datasets as already approved inputs.

\subsubsection{Critical analysis of the research item}
Huang et al. are listed in the proposal as temporal EHR work, connecting deep learning and multimodal models with clinical note comprehension.\cite{zeestB:huang2019} The relevance for Zeest is the emphasis on temporal or multimodal clinical understanding, which aligns with the project's longitudinal memory question. The limitation is that the proposal citation alone does not establish a final Zeest model, dataset or metric. The report should use such work as background for why temporal context matters, while keeping the actual experimental protocol explicit and planned. Across the dataset literature, the common lesson is that reproducible evaluation depends on case definition, permitted access, preprocessing, reference judgement and controlled comparisons. Zeest's planned evaluation therefore uses common case inputs across a single-LLM baseline, a memory-augmented variant and the full system. Any future results must come from those executed comparisons, not from claims imported from prior datasets or unrelated benchmark papers.

\subsubsection{Relationship to the proposed research work}
The dataset literature is related to Zeest because the project needs longitudinal cases, possible imaging inputs and safe development data. MIMIC-IV, MIMIC-CXR, CheXpert and Synthea provide candidate directions, but each has access, preprocessing and domain constraints. Zeest's proposed methodology therefore keeps a dataset and model register instead of claiming that every proposal-listed resource is already selected. These resources will support the planned ablation study only after the final domain, permitted use and reference judgement process are confirmed.

\section{Literature Review Summary Table}
Table~\ref{tab:fatima:literature-summary} summarises the principal works reviewed for Zeest. The table mixes research papers, datasets and established clinical systems because the project is an R\&D integration effort. Each item is included only for the capability it supports: clinical reference, record organisation, medical language modelling, imaging, retrieval, orchestration or evaluation data. The limitations column explains why the item does not already solve the Zeest problem. This distinction is necessary because the proposal combines several mature areas into one physician-controlled workflow. No row is treated as proof that Zeest has been implemented or evaluated. Instead, the table identifies the evidence base used to design the planned methodology and to justify the gaps addressed by the proposed system.

\begin{table}[htbp]
\centering
\scriptsize
\caption{Summary of reviewed systems, models, datasets and methods relevant to Zeest.}
\label{tab:fatima:literature-summary}
\begin{tabular}{|p{0.18\textwidth}|p{0.22\textwidth}|p{0.27\textwidth}|p{0.20\textwidth}|}
\hline
Work & Method or capability & Relevance to Zeest & Limitation for Zeest \\ \hline
Isabel DDx \cite{zeestB:isabel} & Differential diagnosis support system. & Shows clinical value of diagnosis assistance. & Not a longitudinal multimodal agent workflow. \\ \hline
UpToDate \cite{zeestB:uptodate} & Evidence-based clinical reference. & Motivates guideline and citation access. & Does not organise patient memory or approve drafts. \\ \hline
OpenMRS \cite{zeestB:openmrs} & Open source EMR platform. & Supports patient and encounter record thinking. & Not a reasoning, critic and RAG system. \\ \hline
Med-PaLM 2 \cite{zeestB:medpalm2} & Medical LLM alignment and evaluation. & Supports medical-language reasoning potential. & Benchmark QA is not patient-level workflow validation. \\ \hline
GPT-4 medical challenges \cite{zeestA:noriMedical} & Medical examination and benchmark analysis. & Motivates powerful LLM reasoning with caution. & Does not provide persistent memory or approval control. \\ \hline
ClinicalBERT \cite{zeestB:clinicalbert} & Clinical text embeddings. & Supports clinical semantic representation. & Not a complete decision-support workflow. \\ \hline
BioGPT \cite{zeestB:biogpt} & Biomedical generative pre-training. & Supports biomedical text generation and mining. & Does not ground drafts in patient evidence. \\ \hline
CheXNet \cite{zeestB:chexnet} & Chest X-ray pneumonia classifier. & Candidate style for existing imaging model use. & Narrow imaging task without longitudinal context. \\ \hline
CheXpert \cite{zeestB:chexpert} & Chest radiograph dataset and benchmark. & Candidate data/model domain for imaging support. & Access and final domain selection remain unresolved. \\ \hline
BioViL-T \cite{zeestB:biovilt} & Biomedical vision-language temporal learning. & Supports multimodal and temporal motivation. & Does not define Zeest's physician approval workflow. \\ \hline
ReAct \cite{zeestB:react} & Reasoning and acting with tools. & Motivates stepwise agent interaction. & Not a clinical safety or memory architecture. \\ \hline
RAG \cite{zeestA:rag} & Retrieval-augmented generation. & Foundation for cited evidence retrieval. & General method does not ensure clinical sufficiency. \\ \hline
MedRAG \cite{zeestB:medrag} & Medical retrieval-augmented generation benchmark. & Supports medical evidence retrieval evaluation. & Medical QA differs from encounter review. \\ \hline
MedCPT \cite{zeestB:medcpt} & Biomedical retrieval pre-training. & Supports domain-aware evidence search. & Does not select Zeest's final retriever or corpus. \\ \hline
MIMIC-IV \cite{zeestA:mimicIV} & De-identified EHR dataset. & Candidate source for longitudinal cases. & Access, preprocessing and use conditions must be confirmed. \\ \hline
MIMIC-CXR \cite{zeestB:mimiccxr} & De-identified chest radiographs and reports. & Candidate source for imaging and notes. & Does not by itself provide a complete evaluation protocol. \\ \hline
Synthea \cite{zeestB:synthea} & Synthetic patient generation. & Supports safe development case creation. & Synthetic cases may not capture real clinical complexity. \\ \hline
\end{tabular}
\end{table}

\section{Conclusion}
The reviewed literature supports the need for Zeest but also clarifies its boundaries. Clinical support systems and EMR platforms show the importance of diagnosis assistance, evidence access and patient records, yet they do not combine these functions into a longitudinal agentic workflow. Medical LLMs and biomedical language models provide useful reasoning and representation capabilities, but their benchmark performance does not remove the need for evidence, critic review and physician approval. Imaging and vision-language work supports the integration of existing validated models, while also showing why imaging findings must remain tied to their source and task. RAG and medical retrieval research motivate cited evidence retrieval, but retrieval quality must be tested rather than assumed. Dataset work provides possible resources for planned evaluation, subject to access and domain constraints. The resulting gap is the controlled integration of patient memory, multimodal findings, evidence retrieval, critic review and physician decision logging. Chapter~5 uses this gap to define Zeest's proposed methodology and ablation protocol.

\endgroup
